---
tipo: Múltipla Escolha
dificuldade: Média
nivel: médio
disciplina: Matemática
assunto: Juros Simples
gabarito: E
resposta:
tags: [porcentagem, imposto de renda, rendimento]
fonte: concurso
concurso:
banca: Cesgranrio
ano: 1995
numero: 1
cargo:
prova:
---
\question
Os rendimentos das cadernetas de poupanca sao isentos de imposto de renda, os dos fundos de commodities e os dos fundos de renda fixa sao tributados em 25\% e 30\%, respectivamente, da valorizacao que exceder a variacao da UFIR. Suponhamos que para um proximo mes, as previsoes sejam que a UFIR aumente 1,8\% e que as cadernetas, os fundos de commodities e os de renda fixa rendam 2,2\%, 2,6\% e 2,8\%, respectivamente, antes do desconto do Imposto de Renda. Se as previsoes se confirmassem, a melhor e a pior das aplicacoes seriam, respectivamente:
\begin{choices}
\choice poupanca e commodities.
\choice commodities e renda fixa.
\choice commodities e poupanca.
\choice renda fixa e commodities.
\CorrectChoice renda fixa e poupanca.
\end{choices}

---
tipo: Múltipla Escolha
dificuldade: Média
nivel: médio
disciplina: Matemática
assunto: Vetores
gabarito: C
resposta:
tags: [produto escalar, norma]
fonte: concurso
concurso:
banca: Cesgranrio
ano: 1995
numero: 2
cargo:
prova:
---
\question
Se $\lVert \vec{u} + \vec{v} \rVert = 7$ e $\lVert \vec{u} - \vec{v} \rVert = 5$, o valor do produto escalar $\vec{u} \cdot \vec{v}$ e:
\begin{choices}
\choice $8$
\choice $7$
\CorrectChoice $6$
\choice $5$
\choice $4$
\end{choices}

---
tipo: Múltipla Escolha
dificuldade: Média
nivel: médio
disciplina: Matemática
assunto: Porcentagem
gabarito: E
resposta:
tags: [variacao percentual, cambio]
fonte: concurso
concurso:
banca: Cesgranrio
ano: 1995
numero: 3
cargo:
prova:
---
\question
Em 6 de setembro de 1994, os jornais noticiavam que uma grande empresa havia convertido seus precos para reais usando $1$ real $= 2400$ cruzeiros reais e nao $1$ real $= 2750$ cruzeiros reais. Ao fazer isso, nessa empresa, os precos:
\begin{choices}
\choice baixaram cerca de $12,7\%$.
\choice baixaram cerca de $14,6\%$.
\choice aumentaram cerca de $12,7\%$.
\choice aumentaram cerca de $13,2\%$.
\CorrectChoice aumentaram cerca de $14,6\%$.
\end{choices}

---
tipo: Múltipla Escolha
dificuldade: Média
nivel: médio
disciplina: Matemática
assunto: Geometria Espacial
gabarito: A
resposta:
tags: [cilindro, volume, proporcionalidade]
fonte: concurso
concurso:
banca: Cesgranrio
ano: 1995
numero: 4
cargo:
prova:
---
\question
Um salame tem a forma de um cilindro reto com $40$ cm de altura e pesa $1$ kg. Tentando servir um fregues que queria meio quilo de salame, Joao cortou um pedaco, obliquamente, de modo que a altura do pedaco varia entre $22$ cm e $26$ cm. O peso do pedaco e de:
\begin{choices}
\CorrectChoice $600$ g
\choice $610$ g
\choice $620$ g
\choice $630$ g
\choice $640$ g
\end{choices}

---
tipo: Múltipla Escolha
dificuldade: Média
nivel: médio
disciplina: Matemática
assunto: Geometria Espacial
gabarito: A
resposta:
tags: [poliedros, formula de Euler]
fonte: concurso
concurso:
banca: Cesgranrio
ano: 1995
numero: 5
cargo:
prova:
---
\question
Um poliedro convexo tem $14$ vertices. Em $6$ desses vertices concorrem $4$ arestas, em $4$ desses vertices concorrem $3$ arestas e, nos demais vertices, concorrem $5$ arestas. O numero de faces desse poliedro e igual a:
\begin{choices}
\CorrectChoice $16$
\choice $18$
\choice $24$
\choice $30$
\choice $44$
\end{choices}

---
tipo: Múltipla Escolha
dificuldade: Média
nivel: médio
disciplina: Matemática
assunto: Trigonometria
gabarito: C
resposta:
tags: [identidades trigonometricas]
fonte: concurso
concurso:
banca: Cesgranrio
ano: 1995
numero: 6
cargo:
prova:
---
\question
Se $\sen x - \cos x = \frac{1}{2}$, o valor de $\sen x \cos x$ e igual a:
\begin{choices}
\choice $-\frac{3}{16}$
\choice $-\frac{3}{8}$
\CorrectChoice $\frac{3}{8}$
\choice $\frac{3}{4}$
\choice $\frac{3}{2}$
\end{choices}

---
tipo: Múltipla Escolha
dificuldade: Fácil
nivel: médio
disciplina: Matemática
assunto: Razões Trigonométricas no Triângulo Retângulo
gabarito: B
resposta:
tags: [escada, seno]
fonte: concurso
concurso:
banca: Cesgranrio
ano: 1995
numero: 7
cargo:
prova:
---
\question
Uma escada de $2$ m de comprimento esta apoiada no chao e em uma parede vertical. Se a escada faz $30^\circ$ com a horizontal, a distancia do topo da escada ao chao e de:
\begin{choices}
\choice $0,5$ m
\CorrectChoice $1$ m
\choice $1,5$ m
\choice $1,7$ m
\choice $2$ m
\end{choices}

---
tipo: Múltipla Escolha
dificuldade: Média
nivel: médio
disciplina: Matemática
assunto: Números Complexos
gabarito: B
resposta:
tags: [lugar geometrico]
fonte: concurso
concurso:
banca: Cesgranrio
ano: 1995
numero: 8
cargo:
prova:
---
\question
O lugar geometrico das imagens dos complexos $z$, tais que $z^2$ e real, e:
\begin{choices}
\choice um par de retas paralelas.
\CorrectChoice um par de retas concorrentes.
\choice uma reta.
\choice uma circunferencia.
\choice uma parabola.
\end{choices}

---
tipo: Múltipla Escolha
dificuldade: Média
nivel: médio
disciplina: Matemática
assunto: Análise Combinatória
gabarito: D
resposta:
tags: [arranjos, contagem]
fonte: concurso
concurso:
banca: Cesgranrio
ano: 1995
numero: 9
cargo:
prova:
---
\question
Durante a Copa do Mundo, que foi disputada por $24$ paises, as tampinhas de Coca-Cola traziam palpites sobre os paises que se classificariam nos tres primeiros lugares (por exemplo: $1^\circ$ lugar, Brasil; $2^\circ$ lugar, Nigeria; $3^\circ$ lugar, Holanda).

Se, em cada tampinha, os tres paises sao distintos, quantas tampinhas diferentes poderiam existir?
\begin{choices}
\choice $69$
\choice $2024$
\choice $9562$
\CorrectChoice $12144$
\choice $13824$
\end{choices}

---
tipo: Múltipla Escolha
dificuldade: Média
nivel: médio
disciplina: Matemática
assunto: Trigonometria
gabarito: C
resposta:
tags: [rumo, navegacao, vetor]
fonte: concurso
concurso:
banca: Cesgranrio
ano: 1995
numero: 10
cargo:
prova:
---
\question
Um navegador devia viajar durante duas horas, no rumo nordeste, para chegar a certa ilha. Enganou-se, e navegou duas horas no rumo norte. Tomando, a partir dai, o rumo correto, em quanto tempo, aproximadamente, chegara a ilha?
\begin{choices}
\choice $30$ min.
\choice $1$ h.
\CorrectChoice $1$ h $30$ min.
\choice $2$ h.
\choice $2$ h $15$ min.
\end{choices}

---
tipo: Múltipla Escolha
dificuldade: Média
nivel: médio
disciplina: Matemática
assunto: Geometria Analítica
gabarito: B
resposta:
tags: [equacao da reta, figura]
fonte: concurso
concurso:
banca: Cesgranrio
ano: 1995
numero: 11
cargo:
prova:
---
\question
A equacao da reta mostrada na figura a seguir e:
\begin{choices}
\choice $3x + 4y - 12 = 0$
\CorrectChoice $3x - 4y + 12 = 0$
\choice $4x + 3y + 12 = 0$
\choice $4x - 3y - 12 = 0$
\choice $4x - 3y + 12 = 0$
\end{choices}

% figura presente no enunciado original: reta interceptando os eixos em (-4,0) e (0,3)

---
tipo: Múltipla Escolha
dificuldade: Média
nivel: médio
disciplina: Matemática
assunto: Porcentagem
gabarito: D
resposta:
tags: [funcao quadratica, desconto]
fonte: concurso
concurso:
banca: Cesgranrio
ano: 1995
numero: 12
cargo:
prova:
---
\question
Uma loja esta fazendo uma promocao na venda de balas: "Compre $x$ balas e ganhe $x\%$ de desconto". A promocao e valida para compras de ate $60$ balas, caso em que e concedido o desconto maximo de $60\%$. Alfredo, Beatriz, Carlos e Daniel compraram $10$, $15$, $30$ e $45$ balas, respectivamente. Qual deles poderia ter comprado mais balas e gasto a mesma quantia, se empregasse melhor seus conhecimentos de Matematica?
\begin{choices}
\choice Alfredo.
\choice Beatriz.
\choice Carlos.
\CorrectChoice Daniel.
\choice Nenhum.
\end{choices}

---
tipo: Múltipla Escolha
dificuldade: Média
nivel: médio
disciplina: Matemática
assunto: Função Logarítmica
gabarito: B
resposta:
tags: [logaritmo decimal, propriedades]
fonte: concurso
concurso:
banca: Cesgranrio
ano: 1995
numero: 13
cargo:
prova:
---
\question
Se $\log_{10} 123 = 2,09$, o valor de $\log_{10} 1,23$ e:
\begin{choices}
\choice $0,0209$
\CorrectChoice $0,09$
\choice $0,209$
\choice $1,09$
\choice $1,209$
\end{choices}

---
tipo: Múltipla Escolha
dificuldade: Média
nivel: médio
disciplina: Matemática
assunto: Geometria Espacial
gabarito: E
resposta:
tags: [reta e plano, angulo]
fonte: concurso
concurso:
banca: Cesgranrio
ano: 1995
numero: 14
cargo:
prova:
---
\question
$A$ e um ponto nao pertencente a um plano $P$. O numero de retas que contem $A$ e fazem um angulo de $45^\circ$ com $P$ e igual a:
\begin{choices}
\choice $0$.
\choice $1$.
\choice $2$.
\choice $4$.
\CorrectChoice infinito.
\end{choices}

---
tipo: Múltipla Escolha
dificuldade: Média
nivel: médio
disciplina: Matemática
assunto: Geometria Analítica
gabarito: A
resposta:
tags: [area de triangulo, coordenadas]
fonte: concurso
concurso:
banca: Cesgranrio
ano: 1995
numero: 15
cargo:
prova:
---
\question
A area do triangulo, cujos vertices sao $(1, 2)$, $(3, 4)$ e $(4, -1)$, e igual a:
\begin{choices}
\CorrectChoice $6$.
\choice $8$.
\choice $9$.
\choice $10$.
\choice $12$.
\end{choices}

---
tipo: Múltipla Escolha
dificuldade: Média
nivel: médio
disciplina: Matemática
assunto: Função Polinomial
gabarito: C
resposta:
tags: [grafico, polinomio, figura]
fonte: concurso
concurso:
banca: Cesgranrio
ano: 1995
numero: 16
cargo:
prova:
---
\question
O grafico do polinomio $P(x) = x^3 - x^2$ e:
\begin{choices}
\choice alternativa a da figura original.
\choice alternativa b da figura original.
\CorrectChoice alternativa c da figura original.
\choice alternativa d da figura original.
\choice alternativa e da figura original.
\end{choices}

% figura presente no enunciado original: cinco graficos rotulados de a) a e)

---
tipo: Múltipla Escolha
dificuldade: Média
nivel: superior
disciplina: Matemática
assunto: Limites
gabarito: E
resposta:
tags: [infinito, funcao racional]
fonte: concurso
concurso:
banca: Cesgranrio
ano: 1995
numero: 17
cargo:
prova:
---
\question
O valor de $\lim_{x \to 0} \frac{1}{x^2}$ e:
\begin{choices}
\choice $-\infty$
\choice $-1$
\choice $0$
\choice $1$
\CorrectChoice $\infty$
\end{choices}

---
tipo: Múltipla Escolha
dificuldade: Fácil
nivel: superior
disciplina: Matemática
assunto: Derivadas
gabarito: D
resposta:
tags: [cinematica, aceleracao]
fonte: concurso
concurso:
banca: Cesgranrio
ano: 1995
numero: 18
cargo:
prova:
---
\question
Uma particula se move sobre o eixo das abscissas, de modo que sua velocidade no instante $t$ segundos e $v = t^2$ metros por segundo.

A aceleracao dessa particula no instante $t = 2$ segundos e, em metros por segundo quadrado, igual a:
\begin{choices}
\choice $1$.
\choice $2$.
\choice $3$.
\CorrectChoice $4$.
\choice $6$.
\end{choices}

---
tipo: Múltipla Escolha
dificuldade: Média
nivel: médio
disciplina: Matemática
assunto: Geometria Plana
gabarito: A
resposta:
tags: [paralelogramo, area, figura]
fonte: concurso
concurso:
banca: Cesgranrio
ano: 1995
numero: 19
cargo:
prova:
---
\question
$ABCD$ e um paralelogramo e $M$ e o ponto medio do lado $AB$. As retas $CM$ e $BD$ dividem o paralelogramo em quatro partes. Se a area do paralelogramo e $24$, as areas de I, II, III e IV sao, respectivamente, iguais a:
\begin{choices}
\CorrectChoice $10$, $8$, $4$ e $2$.
\choice $10$, $9$, $3$ e $2$.
\choice $12$, $6$, $4$ e $2$.
\choice $16$, $4$, $3$ e $1$.
\choice $17$, $4$, $2$ e $1$.
\end{choices}

% figura presente no enunciado original: paralelogramo ABCD com M no lado AB e regioes I, II, III e IV
